\documentclass[10pt,a4paper]{amsart}
\usepackage[margin=19mm]{geometry}
\usepackage{amsmath,amssymb,mathtools}
\usepackage[scr=rsfs,cal=euler]{mathalfa}
\usepackage{xfrac}
\usepackage[unicode,hidelinks,bookmarks=false]{hyperref}
\newcommand{\ie}{i.e.\ }
\newcommand{\cf}{cf.\ }
\newcommand{\resp}{respectively\ }
\newcommand{\Subsection}{\subsection}
\providecommand{\nobreakdash}{\nobreak\hbox{-}}
\setlength{\emergencystretch}{2em}
\multlinegap0pt
\usepackage{amssymb}
\usepackage{enumerate}
\usepackage[shortlabels]{enumitem}
\usepackage{xcolor}
\newtheorem{proposition}{Proposition}
\newcommand{\bL}{\mathbb{L}}
\newcommand{\bR}{\mathbb{R}}
\newcommand{\definedas}{\mathrel{\raise.095ex\hbox{\rm :}\mkern-5.2mu=}}
\renewcommand{\phi}{\varphi}
\newcommand{\scal}{\mathrm{Scal}}
\DeclareMathOperator{\vol}{Vol}
\title{A new approach towards the construction of initial data in general relativity with positive Yamabe invariant and arbitrary mean curvature}
\author{Armand Coudray, Romain Gicquaud}
\date{Source: arXiv:2603.21214v1 (CC BY 4.0)}
\begin{document}
\maketitle

\noindent\emph{Source and licence.} A new approach towards the construction of initial data in general relativity with positive Yamabe invariant and arbitrary mean curvature. Version arXiv:2603.21214v1 (CC BY 4.0), distributed by arXiv under the Creative Commons Attribution 4.0 International licence, http://creativecommons.org/licenses/by/4.0/, as recorded in the arXiv licence metadata for this exact version and shown on the abstract page https://arxiv.org/abs/2603.21214v1. Full licence notice: \texttt{licenses/arxiv-2603.21214-CC-BY-4.0.txt}.\par\smallskip

\noindent\textit{Editorial scope.} Counted unit: the article's proposition propBootstrap (source0.tex lines 1263--1276). The article's complete original proof of it is delivered with it (source0.tex lines 1278--1387, one \texttt{proof} environment of 110 lines). No mathematics is written, repaired or re-derived: the statement and the proof are the article's own lines, in the article's own notation, and no retained line is substituted or re-flowed.  Retained uncounted as the source-local context the proof needs: lines 680--682; lines 693--706; lines 1006--1011; lines 1209--1216.  Nothing was omitted from any retained span, so the package carries no omission record.  No retained line contains a citation, so the wrapper carries no bibliography and no bibliography entry.  No retained line contains a figure environment, an image-inclusion command or drawing code, so no image is delivered and the package declares no asset dependency.  Editorial formatting only, in addition: an AMS \texttt{amsart} preamble that restates the article's own declarations and macros used by the retained lines, namely the environments \texttt{proposition} on separate counters, together with the standard packages those lines need.  The retained environments print in this document as Proposition 1 in order of appearance, whereas the article prints the same items with the numbering of its own full document.  The complete native source of the article as distributed in the arXiv e-print for this exact version is bundled unchanged as \texttt{sources/196926/source0.tex}.
\begin{equation} \label{eqLich}
    - \frac{4(n-1)}{n-2} \Delta \phi + \scal\,\phi + \frac{n-1}{n} \tau^2 \phi^{N-1} = \frac{A^2}{\phi^{N+1}}.
\end{equation}

\begin{proposition}\label{propEstimateLichUpper}
    Let $\phi$ be the solution to Equation~\eqref{eqLich}, with $A \in L^q(M,
        \bR)$, $q \geq 2$, $A \not\equiv 0$. Then,

    \begin{itemize}
        \item If $q < n$, we have $\left\|\phi\right\|_{L^r}^{N+2} \lesssim \|A\|_{L^q}^2$
              with $r$ such that
              \[
                  \frac{1}{q} = \frac{1}{n} + \frac{2(n-1)}{n-2} \frac{1}{r}.
              \]
        \item If $q \geq n$, we have $\left\|\phi\right\|_{L^r}^{N+2} \lesssim \|A\|_{L^q}^2$
              for any $r \geq N$.
    \end{itemize}
\end{proposition}

\begin{proposition}\label{propVector}
    Assume that $(M, g)$ has no non-trivial conformal Killing vector field, i.e.\
    no non-zero vector field $V$ such that $\bL V \equiv 0$. Then the operator
    $\Delta_{\bL}: W^{2, q}(M, TM) \to L^q(M, TM)$ is an isomorphism for all $q \in
        (1, p]$.
\end{proposition}

\begin{proposition}\label{propStability}
	Let $q = \frac{2n(n-1)}{3n-2}$ be such that $\frac{1}{q} = \frac{1}{n} +
		\frac{1}{2(n-1)}$. If
	\[
		R \leq \left(\frac{\|d\tau\|_{L^n}}{\|d\tau\|_{L^q}}\right)^n V_{\max},
	\]
	then the set $\Omega_0$ is stable under $\Phi$.
\end{proposition}

\begin{proposition}\label{propBootstrap}
	Assume that $\|\sigma\|_{L^{2p}} \leq 1$. There exist an integer $K \geq 0$
	depending only on $n$, and a constant $C > 0$ depending on $(M, g)$, $p$,
	$V_{\max}$, $\|d\tau\|_{L^n}$, and $\|\sigma\|_{L^{2p}}$, such that for all
	$\phi \in \Phi^K\left(\Omega_0\right)$, where $\Phi^K$ is the $K$-th iterate of $\Phi$,
	\[
		\left\|\phi^N\right\|_{L^{2p}} \leq C \|\sigma\|_{L^{2p}}^{\frac{n}{n-1}}.
	\]
	In particular, there exists a constant $C' > 0$ such that
	\[
		\|\bL W\|_{L^{2p}} \leq C' \|\sigma\|_{L^{2p}}^{\frac{n}{n-1}}
	\]
	for all $W = \mathrm{Vect}(\phi)$, with $\phi \in \Phi^K(\Omega_0)$.
\end{proposition}

\begin{proof}
	By Proposition~\ref{propStability}, $\Phi(\Omega_0) \subset \Omega_0$. We
	construct a decreasing sequence of closed subsets $\Omega_k \subset \Omega_0$
	and an increasing sequence of exponents $p_0 < p_1 < \cdots$ defined by
	\[
		p_0 \definedas \frac{N}{2}+1,
		\qquad
		\frac{1}{p_{k+1}} - 1 = \frac{n}{n-1} \left(\frac{1}{p_k} - 1\right),
	\]
	so that
	\begin{equation}\label{eqBootstrap}
		\frac{1}{p_k} = 1 - \left(\frac{n}{n-1}\right)^k \left(1 - \frac{1}{p_0}\right)
		= 1 - \left(\frac{n}{n-1}\right)^k \frac{n}{2(n-1)}.
	\end{equation}

	From formula~\eqref{eqBootstrap}, $1/p_k \to -\infty$, so there exists a
	smallest integer $K_0 \geq 0$ such that $\frac{1}{p_{K_0}} \leq \frac{1}{n}$.
	We claim that $\frac{1}{p_{K_0}} \in \left(0, \frac{1}{n}\right)$. Indeed, for
	the positivity, the recurrence gives, for any $k < K_0$ (i.e.\
	$\frac{1}{p_k} > \frac{1}{n}$),
	\[
		\frac{1}{p_{k+1}} = 1 + \frac{n}{n-1}\left(\frac{1}{p_k} - 1\right)
		> 1 + \frac{n}{n-1}\left(\frac{1}{n} - 1\right) = 0,
	\]
	so in particular $\frac{1}{p_{K_0}} > 0$. For the strict upper bound, note that
	$\frac{1}{p_{K_0}} \in \left\{0, \frac{1}{n}\right\}$ would require, by
	formula~\eqref{eqBootstrap}, that $\left(\frac{n}{n-1}\right)^{K_0+1} = 2$ or
	$\left(\frac{n}{n-1}\right)^{K_0+2} = 2$. Both are impossible since
	$n/(n-1)$ is rational whereas $2^{1/m}$ is irrational for every $m \geq 2$.
	Hence $\frac{1}{p_{K_0}} \in \left(0, \frac{1}{n}\right)$, i.e.\ $p_{K_0} > n$.

	We define $\Omega_k$ by associating to each exponent $p_k$ a radius $R_k
		\definedas C_k \|\sigma\|_{L^{2p}}^{\frac{n}{n-1}}$, where $C_k$ is a positive
	constant to be specified inductively:
	\[
		\Omega_k \definedas \left\{\phi \in L^r(M, \bR_+) \;\middle|\;
		\|\phi^N\|_{L^{p_j}} \leq R_j\ \text{for}\ 0 \leq j \leq k\right\}.
	\]
	For the base case $k = 0$, we set $R_0$ by
	\[
		R_0^{2\frac{n-1}{n}} \definedas \frac{1}{\delta} \|\sigma\|_{L^{2p}}^2 \vol(M, g)^{1-\frac{1}{p}}
		\geq \frac{1}{\delta} \|\sigma\|_{L^2}^2,
	\]
	where the last inequality is H\"older's. By definition of $\Omega_0$, any $\phi
		\in \Omega_0$ satisfies $\|\phi^N\|_{L^{p_0}} \leq R_0$.

	\medskip\noindent\textit{Inductive step ($p_k < n$).}
	Let $\phi \in \Omega_k$ and $W = \mathrm{Vect}(\phi)$. Setting $q_k \definedas
		\left(\frac{1}{p_k}+\frac{1}{n}\right)^{-1}$ and applying elliptic regularity
	(Proposition~\ref{propVector}) followed by the Sobolev embedding theorem gives
	\begin{equation}\label{eqEstimateVector}
		\|\bL W\|_{L^{p_k}} \lesssim \|\phi^N\|_{L^{p_k}} \|d\tau\|_{L^n} \leq R_k \|d\tau\|_{L^n}.
	\end{equation}
	Since $p_k < n$, Proposition~\ref{propEstimateLichUpper} applied to $\phi' =
		\Phi(\phi)$ gives
	\[
		\|(\phi')^N\|_{L^{p_{k+1}}}^{2\frac{n-1}{n}}
		\lesssim \|A\|_{L^{p_k}}^2
		\lesssim \|\sigma\|_{L^{p_k}}^2 + \|\bL W\|_{L^{p_k}}^2
		\lesssim \|\sigma\|_{L^{2p}}^2 + R_k^2,
	\]
	where the exponents $p_k$ and $p_{k+1}$ satisfy $\frac{1}{p_k} = \frac{1}{n} +
		\frac{n-1}{n} \frac{1}{p_{k+1}}$, which is equivalent to the recurrence
	defining $(p_k)$. Let $\Lambda_{k+1}^2$ denote the implicit constant:
	\[
		\|(\phi')^N\|_{L^{p_{k+1}}}^{2\frac{n-1}{n}}
		\leq \Lambda_{k+1}^2 \left(\|\sigma\|_{L^{2p}}^2 + R_k^2\right).
	\]
	Since $R_k = C_k \|\sigma\|_{L^{2p}}^{\frac{n}{n-1}}$ and $\|\sigma\|_{L^{2p}}
		\leq 1$, we get
	\begin{align*}
		\|(\phi')^N\|_{L^{p_{k+1}}}^{2\frac{n-1}{n}}
		 & \leq \Lambda_{k+1}^2 \|\sigma\|_{L^{2p}}^2 \left(1 + C_k^2 \|\sigma\|_{L^{2p}}^{\frac{2}{n-1}}\right)
		\leq \Lambda_{k+1}^2 \left(1 + C_k^2\right) \|\sigma\|_{L^{2p}}^2.
	\end{align*}
	Setting $C_{k+1}$ by $C_{k+1}^{2\frac{n}{n-1}} \definedas \Lambda_{k+1}^2 (1 + C_k^2)$, we
	obtain $\|(\phi')^N\|_{L^{p_{k+1}}} \leq R_{k+1}$, hence $\Phi(\Omega_k) \subset
		\Omega_{k+1}$ and by induction $\Phi^k(\Omega_0) \subset \Omega_k$ for all
	$0 \leq k \leq K_0$.

	\medskip\noindent\textit{Terminal step ($p_{K_0} \geq n$).}
	We apply Proposition~\ref{propEstimateLichUpper} (second case) to
	$\phi \in \Omega_{K_0}$: since $p_{K_0} \geq n$, for any $r \geq 1$ we have
	\[
		\|(\Phi(\phi))^N\|_{L^r}^{2\frac{n-1}{n}}
		\lesssim \|A\|_{L^{p_{K_0}}}^2
		\lesssim \|\sigma\|_{L^{2p}}^2 + R_{K_0}^2
		\leq (1 + C_{K_0}^2)\|\sigma\|_{L^{2p}}^2,
	\]
	where the last inequality uses $\|\sigma\|_{L^{2p}} \leq 1$. Since $2(n-1)/n =
		1 + 2/N$, we conclude that
	\[
		\|(\Phi(\phi))^N\|_{L^r} \leq C_r \|\sigma\|_{L^{2p}}^{\frac{n}{n-1}},
	\]
	where $C_r$ depends on $(M, g)$, $p$, $K_0$, and $r$, but not on $\phi$.
	Setting $K \definedas K_0 + 1$, we have shown that every $\phi \in
		\Phi^K(\Omega_0) \subset \Phi(\Omega_{K_0})$ satisfies the claimed bound on
	$\|\phi^N\|_{L^r}$. Taking $r = 2p$ gives the first part of the proposition.

	\medskip

	For the second part, let $\phi \in \Phi^K(\Omega_0)$ and $W =
		\mathrm{Vect}(\phi)$. Applying the estimate~\eqref{eqEstimateVector} with $p_k$
	replaced by $2p$ and using the bound just established yields
	\[
		\|\bL W\|_{L^{2p}} \lesssim \|\phi^N\|_{L^{2p}} \|d\tau\|_{L^n}
		\lesssim \|\sigma\|_{L^{2p}}^{\frac{n}{n-1}} \|d\tau\|_{L^n},
	\]
	which completes the proof.
\end{proof}

\end{document}
